% Options for packages loaded elsewhere
\PassOptionsToPackage{unicode}{hyperref}
\PassOptionsToPackage{hyphens}{url}
\documentclass[
  ignorenonframetext,
]{beamer}
\newif\ifbibliography
\usepackage{pgfpages}
\setbeamertemplate{caption}[numbered]
\setbeamertemplate{caption label separator}{: }
\setbeamercolor{caption name}{fg=normal text.fg}
\beamertemplatenavigationsymbolsempty
% remove section numbering
\setbeamertemplate{part page}{
  \centering
  \begin{beamercolorbox}[sep=16pt,center]{part title}
    \usebeamerfont{part title}\insertpart\par
  \end{beamercolorbox}
}
\setbeamertemplate{section page}{
  \centering
  \begin{beamercolorbox}[sep=12pt,center]{section title}
    \usebeamerfont{section title}\insertsection\par
  \end{beamercolorbox}
}
\setbeamertemplate{subsection page}{
  \centering
  \begin{beamercolorbox}[sep=8pt,center]{subsection title}
    \usebeamerfont{subsection title}\insertsubsection\par
  \end{beamercolorbox}
}
% Prevent slide breaks in the middle of a paragraph
\widowpenalties 1 10000
\raggedbottom
\AtBeginPart{
  \frame{\partpage}
}
\AtBeginSection{
  \ifbibliography
  \else
    \frame{\sectionpage}
  \fi
}
\AtBeginSubsection{
  \frame{\subsectionpage}
}
\usepackage{iftex}
\ifPDFTeX
  \usepackage[T1]{fontenc}
  \usepackage[utf8]{inputenc}
  \usepackage{textcomp} % provide euro and other symbols
\else % if luatex or xetex
  \usepackage{unicode-math} % this also loads fontspec
  \defaultfontfeatures{Scale=MatchLowercase}
  \defaultfontfeatures[\rmfamily]{Ligatures=TeX,Scale=1}
\fi
\usepackage{lmodern}
\ifPDFTeX\else
  % xetex/luatex font selection
\fi
% Use upquote if available, for straight quotes in verbatim environments
\IfFileExists{upquote.sty}{\usepackage{upquote}}{}
\IfFileExists{microtype.sty}{% use microtype if available
  \usepackage[]{microtype}
  \UseMicrotypeSet[protrusion]{basicmath} % disable protrusion for tt fonts
}{}
\makeatletter
\@ifundefined{KOMAClassName}{% if non-KOMA class
  \IfFileExists{parskip.sty}{%
    \usepackage{parskip}
  }{% else
    \setlength{\parindent}{0pt}
    \setlength{\parskip}{6pt plus 2pt minus 1pt}}
}{% if KOMA class
  \KOMAoptions{parskip=half}}
\makeatother
\setlength{\emergencystretch}{3em} % prevent overfull lines
\providecommand{\tightlist}{%
  \setlength{\itemsep}{0pt}\setlength{\parskip}{0pt}}
% Slide layout defaults for warning-clean scientific decks.
\usepackage{etoolbox}
\IfFileExists{xurl.sty}{\usepackage{xurl}}{}
\IfFileExists{seqsplit.sty}{\usepackage{seqsplit}}{\newcommand{\seqsplit}[1]{#1}}
\protected\def\breakseq#1{\seqsplit{#1}}
\protected\def\breaktt#1{\begingroup\ttfamily\seqsplit{#1}\endgroup}
\setlength{\emergencystretch}{6em}
\tolerance=5000
\hbadness=10000
\hfuzz=1pt
\setlength{\tabcolsep}{2pt}
\AtBeginEnvironment{longtable}{\tiny\renewcommand{\arraystretch}{0.86}\setlength{\tabcolsep}{1pt}}
\AtBeginEnvironment{tabular}{\tiny\renewcommand{\arraystretch}{0.86}\setlength{\tabcolsep}{1pt}}

% Natbib and cross-reference fallbacks — slides don't load natbib
% or cleveref, but combined-PDF manuscript prose may emit these
% commands. The fallback renders citations as a bracketed key list
% and cross-references as detokenized labels so slides don't fail on
% undefined control sequences or raw underscores. \providecommand is
% a no-op if packages load later.
\providecommand{\citep}[1]{[#1]}
\providecommand{\citet}[1]{#1}
\providecommand{\citealp}[1]{#1}
\providecommand{\citeauthor}[1]{#1}
\providecommand{\citeyear}[1]{#1}
\providecommand{\cref}[1]{\texttt{\detokenize{#1}}}
\providecommand{\Cref}[1]{\texttt{\detokenize{#1}}}
\usepackage{bookmark}
\IfFileExists{xurl.sty}{\usepackage{xurl}}{} % add URL line breaks if available
\urlstyle{same}
% fallback for those not using the hyperref driver hyperxmp:
\makeatletter
\@ifundefined{xmpquote}{\newcommand{\xmpquote}[1]{#1}}{}
\makeatother
\hypersetup{
  hidelinks,
  pdfcreator={LaTeX via pandoc}}

\author{\texorpdfstring{}{}}
\date{}

\begin{document}

\section{Conclusion}\label{sec:conclusion}

\begin{frame}[fragile,allowframebreaks]{Summary}
\protect\phantomsection\label{summary}
This paper has presented \breaktt{template\_pools\_rules\_tools}, a
meta-project exemplar demonstrating how research software projects
embedded in a monorepo can integrate three categories of shared
resources --- data pools (fonds), governance rules, and executable tools
--- without tight coupling to any specific resource instance. The key
contributions are:

\begin{enumerate}
\item
  \textbf{A four-module architecture} (\texttt{fonds\_reader},
  \texttt{rules\_applier}, \texttt{tools\_invoker},
  \texttt{integration}) that provides a canonical pattern for
  resource-aware project code in the template repository. Each module is
  independently testable and independently deployable.
\item
  \textbf{A typed manifest convention} (\texttt{fonds.yaml},
  \texttt{rules.yaml}, \texttt{tools.yaml}) that makes resource
  capabilities explicit and checkable at pipeline initialisation time,
  shifting failure detection from runtime to startup --- a significant
  improvement for reproducibility \breakseq{[@Wilson2014best]}.
\item
  \textbf{A token injection pipeline} that links manuscript prose to
  integration runtime statistics through \texttt{\{\{integration.*\}\}}
  tokens, ensuring that quantitative claims in the manuscript are always
  generated by the pipeline rather than authored manually. In the
  current run this covered \{\{integration.fonds\_loaded\}\} fonds,
  \{\{integration.rules\_sets\_ok\}\} rule sets,
  \{\{integration.tools\_discovered\}\} tools, and
  \{\{integration.bib\_entries\}\} bibliography entries.
\item
  \textbf{A three-level resilience design} --- resource absence, schema
  malformation, and script absence --- that allows the pipeline to
  degrade gracefully and report failures informatively rather than
  crashing, consistent with best practices for robust research software
  \breakseq{[@Taschuk2017ten]}.
\end{enumerate}
\end{frame}

\begin{frame}[fragile,allowframebreaks]{Design Decisions Revisited}
\protect\phantomsection\label{design-decisions-revisited}
Several design decisions deserve emphasis as lessons for future exemplar
authors:

\textbf{Repo-root-relative discovery} is non-negotiable in a forkable
monorepo. Any hard-coded absolute path or working-directory assumption
will fail when the repository is cloned to a different location. The
\texttt{pathlib.Path(\_\_file\_\_).resolve().parents{[}N{]}} idiom used
throughout \texttt{src/} ensures that discovery works from any working
directory.

\textbf{Graceful fallbacks over exceptions} is the correct trade-off for
a resource-consumption layer. The alternative --- raising
\breaktt{FileNotFoundError} when a fond is missing --- would make the
integration pipeline fragile to the ordering of parallel agent writes.
Returning \texttt{None} and logging a warning allows the pipeline to
produce a complete, if partial, result dict that downstream consumers
can reason about.

\textbf{Real-file tests over mocks} is required for a template exemplar.
Mocks can pass tests while the real discovery logic is silently broken.
Tests that use real file paths exercise the full path-resolution chain
and catch regressions that mocks would miss.

\textbf{Thin orchestration scripts} keep the scripts directory readable
and focus all testable logic in \texttt{src/}. A script that does more
than parse arguments, call a \texttt{src/} function, and write output is
accumulating business logic that belongs elsewhere.
\end{frame}

\begin{frame}[fragile,allowframebreaks]{Future Directions}
\protect\phantomsection\label{future-directions}
The three-layer architecture described here is deliberately extensible.
The most natural extension is a fourth resource category:
\textbf{models}
(\texttt{models/\textless{}scope\textgreater{}/\textless{}name\textgreater{}/})
for pre-trained machine learning models and their inference scripts. The
pattern would follow exactly the same manifest-and-reader design, with
\texttt{models.yaml} declaring type, version, and entrypoints, and a
\texttt{models\_loader} module providing \breaktt{discover\_models()} and
\breaktt{validate\_model\_files\_exist()}. Adding this layer to
\breaktt{template\_pools\_rules\_tools} would require only a new reader
module and a new section in the integration orchestrator --- the
existing architecture places no constraints on the number of resource
categories.

A second direction is \textbf{cross-fond validation}: checking that cite
keys used in a project's \breaktt{manuscript/references.bib} are a subset
of those available in \breaktt{fonds/templates/template\_bibliography/}.
This would require a small cross-reference function in the integration
module and would strengthen the reproducibility guarantee by ensuring
that manuscript citations always trace back to curated, shared sources.

The exemplar is ready to fork. A project team wishing to adopt this
architecture need only copy the four \texttt{src/} modules, update the
resource directory paths in each reader, and replace the exemplar
fond/rules/tool names with their own.
\end{frame}

\end{document}
